\section{Diophantine Equation}
\subsection{Problem 66}
Consider quadratic Diophantine equations of the form:
$x^2 - Dy^2 = 1$
For example, when $D=13$, the minimal solution in $x$ is $649^2 - 13 \times 180^2 = 1$.

It can be assumed that there are no solutions in positive integers when $D$ is square.

By finding minimal solutions in $x$ for $D = \{2, 3, 5, 6, 7\}$, we obtain the following:

\begin{align*}
	3^2 - 2 \times 2^2                        & = 1 \\
	2^2 - 3 \times 1^2                        & = 1 \\
	{\color{red}{\mathbf 9}}^2 - 5 \times 4^2 & = 1 \\
	5^2 - 6 \times 2^2                        & = 1 \\
	8^2 - 7 \times 3^2                        & = 1
\end{align*}

Hence, by considering minimal solutions in
$x$ for $D \le 7$, the largest $x$ is obtained when $D = 5$.

Find the value of $D \le 1000$ in minimal solutions of $x$ for which the largest value of $x$ is obtained.

\subsection{数学解法}
\subsubsection{连分数}
\begin{definition}
	设 \( x_0, x_1, x_2, \dots \) 是一个无穷实数列， \( x_j > 0, j \ge 1 \)。 对缩写的 \( n \ge 0 \)， 我们把表示式
	\begin{equation}\label{eq:cfrac}
		\cfrac{1}{x_0 +
			\cfrac{1}{x_1 +
				\cfrac{1}{x_2 +
					\cfrac{1}{x_3 +
						\cfrac{1}{\ddots +
							\cfrac{1}{x_n}
						}}}}}
	\end{equation}
	称为\textbf{( \( n \) 阶)有限连分数}，它的值是一个实数。当 \( x_0, \dots, x_n
	\)均为整数时称为($n$阶)\em{有限简单连分数}，它的值是一个有理分数。为书写方便，把有限连分数记作
	\begin{equation}
		\langle x_0,x_1,\dots,x_n\rangle.
		\label{eq:ncfrac}
	\end{equation}
	设 \( 0 \le k \le n \) ， 我们把有限连分数
	\begin{equation}
		\langle x_0, x_1, \dots, \x_k\rangle
		\label{eq:kcfrac}
	\end{equation}
	称为是有限连分数\eqref{eq:ncfrac}的第 \( k \)个渐进分数。当\eqref{eq:ncfrac}是有限简单连分数(即 \( x_0,\dots,x_n
	\)均为整数)时，把 \( x_k(0 \le k \le n) \) 称为是它的第 \( k \)
	个部分商。当式\eqref{eq:cfrac}(或\eqref{eq:ncfrac})中的 \( n \to \infty
	\)时，我们把相应的表示式\eqref{eq:cfrac}(或\eqref{eq:cfrac})称为无限连分数，即表示式
	\begin{equation}\label{eq:infcfrac}
		x_0 + \cfrac{1}{x_1 +
				\cfrac{1}{x_2 +
					\cfrac{1}{x_3 +
						\genfrac{}{}{0pt}{}{\vphantom{1}}{\ddots}
					}}},
	\end{equation}
	或者简记为
	\begin{equation}
		\langle x_0, x_1, x_2, \dots\rangle.
		\label{eq:abbrinfcfrac}
	\end{equation}
	当 \( x_j(j \ge 0) \)均为整数时，称\eqref{eq:infcfrac}(或\eqref{eq:abbrinfcfrac})为无限简单连分数。同样的，对任意
	\( k \ge 0 \)，有限连分数\eqref{eq:kcfrac}称为是无限连分数\eqref{eq:infcfrac}(或\eqref{eq:abbrinfcfrac})的第 \( k
	\)个渐近分数；当\eqref{eq:infcfrac}(或\eqref{eq:abbrinfcfrac})是无限简单连分数时， \( x_k(k \ge 0) \)称为是它的第
	\( k \)个部分商。如果存在极限
	\begin{equation}
		\lim_{k\to\infty}\langle x_0,\dots,x_k\rangle = \theta,
		\label{eq:cfraclim}
	\end{equation}
	那么，就说无限连分数\eqref{eq:infcfrac}(或\eqref{eq:abbrinfcfrac})是收敛的， \( \theta \)
	称为是无限连分数\eqref{eq:infcfrac}(或\eqref{eq:abbrinfcfrac})的值，记作
	\begin{equation}
		\langle x_0,x_1,x_2,\dots\rangle = \theta,
		\label{eq:infcfracvalue}
	\end{equation}
	或极限\eqref{eq:cfraclim}不存在，则说无限连分数\eqref{eq:infcfrac}(或\eqref{eq:abbrinfcfrac})是发散的。
\end{definition}
\begin{theorem}
	设 \( x_0, x_1, x_2, \dots \)是无穷实数列， \( x_j > 0, j \ge 1 \)。那么，
	\begin{enumerate}[label=\rm(\roman*)]
		\item 对任意的整数 \( n \ge 1, r \ge 1 \)有
		      \begin{flalign}
			      \langle x_0,\dots,x_{n-1},x_n,\dots,x_{n+r}\rangle \nonumber &                                                                                        \\
			                                                                   & = \langle x_0,\dots,x_{n-1},\langle x_n,\dots,x_{n+4}\rangle\rangle\nonumber           \\
			                                                                   & = \langle x_0,\dots,x_{n-1}, x_{n} + 1 / \langle x_{n+1},\dots,x_{n+r}\rangle\rangle .
		      \end{flalign}
		      特别地有
		      \begin{equation}
			      \langle x_0, \dots, x_{n-1},x_n,x_{n+1}\rangle=\langle x_0,\dots,x_{n-1},x_n+1/x_{n+1}\rangle.
		      \end{equation}
		\item 对任意实数 \( \eta>0 \)及整数 \( n \ge 1 \)，
		      \begin{gather}
			      \langle x_0,\dots, x_{n-1}, x_n\rangle > \langle x_0,\dots,x_{n-1}, x_{n}+\eta\rangle, 2 \nmid n. \\
			      \langle x_0,\dots, x_{n-1}, x_n\rangle < \langle x_0,\dots,x_{n-1},x_n + \eta\rangle, 2 \mid n.
		      \end{gather}
		\item 记
		      \begin{equation}
			      a^{(n)} = \langle x_0, \dots, x_n \rangle.
		      \end{equation}
		      我们有
		      \begin{gather}
			      a^{(n)} > a^{(n + r)}, 2 \nmid n, r \ge 1, \\
			      a^{(n)} < a^{(n + r)}, 2 \mid n, r \ge 1, \\
			      a^{(1)} > a^{(3)} > a^{(5)} > \dots > a^{2s - 1} > \dots, \\
			      a^{(0)} < a^{(2)} < a^{(4)} < \dots < a^{2s} < \dots, \\
			      a^{(2s-1)} > a^{(2t)}, s \ge 1, t \ge 0 .
		      \end{gather}
	\end{enumerate}
\end{theorem}
